% Partie : logique propositionnelle
% Chapitre : syntaxe et sémantique
% Section : correction

\documentclass[../../../include/open-logic-section]{subfiles}

\begin{document}

\olfileid[fr]{pl}{prp}{snd}

\olsection{Correction de la logique propositionnelle}

\begin{thm}[Correction]
\ollabel{thm:soundness}
Si $\Gamma\Proves !A$, alors
$\Gamma\Entails !A$.
\end{thm}

\begin{proof}
Procédons par récurrence sur la longueur d'une
dérivation finie de~$!A$ à partir de~$\Gamma$.
Les tables de vérité des schémas d'axiomes montrent
que ceux-ci sont des tautologies. Leurs instances
sont donc valides, et par conséquent des conséquences
logiques de~$\Gamma$. Si $!A\in\Gamma$,
toute valuation qui satisfait~$\Gamma$ satisfait
également~$!A$.

Pour le pas d'inférence, supposons que $!A$ soit
obtenue par \emph{modus ponens} à partir de
$!C$ et $!C\lif !A$. Par l'hypothèse de
récurrence, $\Gamma\Entails !C$ et
$\Gamma\Entails !C\lif !A$. La deuxième
assertion de la
\olref[pl][syn][sem]{prop:semanticalfacts}
donne alors $\Gamma\Entails !A$.
\end{proof}

\begin{cor}
Si $\Gamma$ est satisfaisable, alors $\Gamma$
est cohérent. En particulier, la logique
propositionnelle est cohérente.
\end{cor}

\begin{proof}
Si $\Gamma$ était incohérent, il permettrait
de dériver~$\lfalse$. Par correction,
on aurait $\Gamma\Entails\lfalse$, ce qui
contredit l'existence d'une valuation
satisfaisant~$\Gamma$. L'ensemble vide est
satisfaisable par toute valuation propositionnelle ;
le système lui-même est donc cohérent.
\end{proof}

\end{document}
